A questão que nos é posta hoje é um belíssimo exemplo de como propriedades aparentemente
simples, quando impostas universalmente, podem forçar uma estrutura extremamente rígida sobre
um objeto algébrico. O traço, uma função aparentemente modesta — a soma dos elementos da
diagonal —, revela-se aqui um poderoso detetive do comportamento interno da matriz.

Vamos à demonstração, que conectará a noção de traço, autovalores e, em última instância, a
nilpotência, em uma cadeia lógica impecável.

\subsection*{Enunciado do Problema:}

Seja $A$ uma matriz complexa $n \times n$ tal que, para todo inteiro positivo $k$, tem-se traço$(A^k)=0$.
Prove que $A$ é nilpotente.

\subsection*{Demonstração}

Nossa estratégia será demonstrar que a condição imposta sobre os traços das potências de $A$
implica que o único autovalor possível para $A$ é o escalar $0$. Uma vez estabelecido este fato,
invocaremos o resultado do nosso primeiro exercício para concluir a prova.

\textbf{Passo 1: A Relação Fundamental entre Traço e Autovalores}

Recordemos dois resultados fundamentais da teoria espectral.

\textbf{Teorema 1:} O traço de uma matriz complexa $A \in M_n(\mathbb{C})$ é igual à soma de seus autovalores,
contados com suas respectivas multiplicidades algébricas.
Se $\lambda_1, \lambda_2, \ldots, \lambda_n$ são os autovalores de $A$, então:

\[
\mathrm{tr}(A) = \sum_{i=1}^n \lambda_i
\]

\textbf{Teorema 2:} Se $\lambda_1, \lambda_2, \ldots, \lambda_n$ são os autovalores de $A$, então os autovalores da matriz $A^k$ (para $k \in \mathbb{Z}^+$) são $\lambda_1^k, \lambda_2^k, \ldots, \lambda_n^k$.

Combinando estes dois teoremas, obtemos uma expressão para o traço de $A^k$ em termos dos
autovalores de $A$:

\[
\mathrm{tr}(A^k) = \sum_{i=1}^n \lambda_i^k
\]

\textbf{Passo 2: Traduzindo a Hipótese do Problema}

A hipótese do problema é que $\mathrm{tr}(A^k)=0$ para todo inteiro positivo $k \geq 1$. Utilizando a relação que
acabamos de estabelecer, esta hipótese se traduz no seguinte sistema infinito de equações sobre os
autovalores de $A$:

\[
\begin{cases}
k=1: & \lambda_1+\lambda_2+\cdots+\lambda_n=0 \\
k=2: & \lambda_1^2+\lambda_2^2+\cdots+\lambda_n^2=0 \\
k=3: & \lambda_1^3+\lambda_2^3+\cdots+\lambda_n^3=0 \\
\quad & \vdots
\end{cases}
\]

Denominamos a soma das $k$-ésimas potências como $p_k=\sum_{i=1}^n \lambda_i^k$. A nossa hipótese é, portanto,
$p_k=0$ para todo $k \geq 1$. Nosso objetivo é mostrar que esta condição força $\lambda_i=0$ para todo $i$.

\textbf{Passo 3: As Identidades de Newton e os Polinômios Simétricos}

Para provar que todos os autovalores são nulos, faremos uso das célebres \textit{Identidades de Newton},
que relacionam as somas de potências $(p_k)$ com os polinômios simétricos elementares $(e_j)$ nas
variáveis $\lambda_1,\ldots,\lambda_n$.

Os polinômios simétricos elementares são definidos como:

\[
e_1=\sum_i \lambda_i, \quad
e_2=\sum_{i<j} \lambda_i\lambda_j, \quad
e_3=\sum_{i<j<k} \lambda_i\lambda_j\lambda_k, \quad
\ldots, \quad
e_n=\lambda_1\lambda_2\cdots\lambda_n
\]

As Identidades de Newton fornecem uma relação recursiva:

\[
ke_k = \sum_{i=1}^k (-1)^{i-1} e_{k-i} p_i
\]

Vamos "abrir" as primeiras identidades para maior clareza:

\begin{itemize}
\item Para $k=1$: $e_1=p_1$
\item Para $k=2$: $2e_2=e_1p_1-p_2$
\item Para $k=3$: $3e_3=e_2p_1-e_1p_2+p_3$
\item E assim por diante.
\end{itemize}

Agora, aplicamos a nossa hipótese, $p_k=0$ para todo $k \geq 1$:

\begin{itemize}
\item $e_1=p_1=0$.
\item $2e_2=e_1p_1-p_2=(0)(0)-0=0 \;\;\Longrightarrow\;\; e_2=0$.
\item $3e_3=e_2p_1-e_1p_2+p_3=(0)(0)-(0)(0)+0=0 \;\;\Longrightarrow\;\; e_3=0$.
\end{itemize}

Por indução, se supomos que $e_1,e_2,\ldots,e_{k-1}$ são todos nulos, a identidade
$ke_k=\sum_{i=1}^k (-1)^{i-1} e_{k-i}p_i$ se torna:

\[
ke_k=e_{k-1}p_1 - e_{k-2}p_2+\cdots+(-1)^{k-2} e_1p_{k-1}+(-1)^{k-1} e_0 p_k.
\]

Como $e_1,\ldots,e_{k-1}$ são nulos e $p_1,\ldots,p_k$ são nulos, todos os termos do lado direito se anulam.
Logo, $ke_k=0$, o que implica $e_k=0$.

Portanto, demonstramos que $e_j=0$ para todo $j=1,2,\ldots,n$.

\textbf{Passo 4: Conexão com o Polinômio Característico}

Qual a relevância dos polinômios simétricos elementares? Eles são, a menos de um sinal, os
coeficientes do polinômio característico de $A$. O polinômio característico $p_A(t)$ tem os autovalores
$\lambda_1,\ldots,\lambda_n$ como suas raízes. Portanto, ele pode ser escrito como:

\[
p_A(t)=(t-\lambda_1)(t-\lambda_2)\cdots(t-\lambda_n)
\]

Expandindo esta forma fatorada, obtemos:

\[
p_A(t)=t^n-e_1t^{n-1}+e_2t^{n-2}-\cdots+(-1)^n e_n
\]

No Passo 3, provamos que $e_1=e_2=\cdots=e_n=0$. Substituindo estes valores na expressão do
polinômio característico:

\[
p_A(t)=t^n-(0)t^{n-1}+(0)t^{n-2}-\cdots+(-1)^n(0)
\]

\[
p_A(t)=t^n
\]

As raízes do polinômio característico são os autovalores da matriz. As raízes de $p_A(t)=t^n$ são $t=0$
com multiplicidade $n$.
Concluímos, de forma inequívoca, que o único autovalor da matriz $A$ é $\lambda=0$.

\textbf{Passo 5: A Conclusão Final}

Chegamos ao ponto crucial. Nós estabelecemos que a hipótese $\mathrm{tr}(A^k)=0$ para todo $k \geq 1$
implica que a matriz $A$ tem somente o autovalor $0$.

Conforme rigorosamente demonstrado no Exercício 1, "se uma matriz $n\times n$ complexa $A$ tem
somente $0$ como autovalor, então $A$ é nilpotente".

Portanto, por consequência direta do resultado anterior e da nossa dedução, a matriz $A$ é nilpotente.

\textbf{Quod Erat Demonstrandum (Q.E.D.)}